\chapter{1990年广东卷（理）}
\section{单选题}

\begin{problem}
如果直线 \(ax+2y+1=0\) 与直线 \(x+y-2=0\) 互相垂直，那么 \(a\) 的值等于（\quad）

\choices
    {\(1\)}
    {\(-\frac{1}{3}\)}
    {\(-\frac{2}{3}\)}
    {\(-2\)}
\end{problem}

\begin{answer}
D
\end{answer}

\begin{solution}
两直线的斜率分别为 \(m_1=-\frac{a}{2}\) 和 \(m_2=-1\). 由于两直线互相垂直，\(m_1m_2=-1\)，所以
\(\left(-\frac{a}{2}\right)(-1)=-1\)，
解得 \(a=-2\). 因而选 D.
\end{solution}

\begin{problem}
如果直线 \(l\) 是平面 \(\alpha\) 的斜线，那么在平面 \(\alpha\) 内（\quad）

\choices
    {不存在与 \(l\) 平行的直线}
    {不存在与 \(l\) 垂直的直线}
    {与 \(l\) 垂直的直线只有一条}
    {与 \(l\) 平行的直线有无穷多条}
\end{problem}

\begin{answer}
A
\end{answer}

\begin{solution}
设 \(P=l\cap\alpha\). 若平面 \(\alpha\) 内存在直线与 \(l\) 平行，则 \(l\) 将与平面 \(\alpha\) 平行或在平面 \(\alpha\) 内，这与 \(l\) 是平面 \(\alpha\) 的斜线矛盾，所以平面 \(\alpha\) 内不存在与 \(l\) 平行的直线.

设 \(l'\) 是 \(l\) 在平面 \(\alpha\) 内的射影. 在平面 \(\alpha\) 内过 \(P\) 作直线 \(m\perp l'\)，则 \(m\) 与 \(l\) 所成的角为直角. 平面 \(\alpha\) 内有无数条与 \(m\) 平行的直线，它们的方向都与 \(l\) 的方向垂直，因而与 \(l\) 所成的角也都是直角. 所以平面 \(\alpha\) 内与 \(l\) 垂直的直线有无数条，只有选项 A 正确.
\end{solution}

\begin{problem}
某乒乓球队有9名队员，其中2名是种子选手，现要挑选5名队员参加比赛，种子选手都必须在内，那么不同的选法共有（\quad）

\choices
    {\(126\) 种}
    {\(84\) 种}
    {\(35\) 种}
    {\(21\) 种}
\end{problem}

\begin{answer}
C
\end{answer}

\begin{solution}
两名种子选手必须入选，先确定这两人，再从其余 \(9-2=7\) 名队员中选出 \(5-2=3\) 名，所以选法数为
\(\mathrm C_7^3=\frac{7\cdot6\cdot5}{3\cdot2\cdot1}=35\).
因此选 C.
\end{solution}

\begin{problem}
把6个不同的元素排成前后两排，每排3个元素，那么不同的排法共有（\quad）

\choices
    {\(36\) 种}
    {\(120\) 种}
    {\(720\) 种}
    {\(1440\) 种}
\end{problem}

\begin{answer}
C
\end{answer}

\begin{solution}
前后两排的三个位置彼此有区别，六个不同元素只需分别占据这六个位置. 因而排法数为
\(6!=6\cdot5\cdot4\cdot3\cdot2\cdot1=720\).
因此选 C.
\end{solution}

\begin{problem}
一个圆台的母线长是上、下底面半径的等差中项，且侧面积为 \(8\pi \mathrm{~cm}^2\)，那么母线长为（\quad）

\choices
    {\(4\mathrm{~cm}\)}
    {\(2\sqrt{2}\mathrm{~cm}\)}
    {\(2\mathrm{~cm}\)}
    {\(\sqrt{2}\mathrm{~cm}\)}
\end{problem}

\begin{answer}
C
\end{answer}

\begin{solution}
设圆台上、下底面半径分别为 \(r\)、\(R\)，母线长为 \(l\). 由题意，\(l\) 是 \(r\)、\(R\) 的等差中项，故
\(l=\frac{r+R}{2}\). 圆台的侧面积为 \(\pi(r+R)l\)，所以 \(\pi(r+R)l=8\pi\).
代入 \(r+R=2l\)，得 \(2\pi l^2=8\pi\)，从而 \(l^2=4\). 因为母线长为正，\(l=2\mathrm{~cm}\)，选 C.
\end{solution}

\begin{problem}
如果复数 \(z=3+a\i\) 满足条件 \(\left| z-2 \right|<2\)，那么实数 \(a\) 的取值范围是（\quad）

\choices
    {\((-2\sqrt{2}, 2\sqrt{2})\)}
    {\((-2,2)\)}
    {\((-1,1)\)}
    {\((-\sqrt{3}, \sqrt{3})\)}
\end{problem}

\begin{answer}
D
\end{answer}

\begin{solution}
由 \(z=3+a\i\)，有 \(z-2=1+a\i\)，于是
\(|z-2|=|1+a\i|=\sqrt{1+a^2}\).
题设不等式等价于 \(\sqrt{1+a^2}<2\)，即 \(a^2<3\). 因此
\(-\sqrt{3}<a<\sqrt{3}\)，
故选 D.
\end{solution}

\begin{problem}
一动点在圆 \(x^2+y^2=1\) 上移动时，它与定点 \((3,0)\) 连线的中点的轨迹方程是（\quad）

\choices
    {\((x+3)^2+y^2=4\)}
    {\((x-3)^2+y^2=1\)}
    {\((2x-3)^2+4y^2=1\)}
    {\(\left( x+\frac{3}{2} \right)^2+y^2=\frac{1}{2}\)}
\end{problem}

\begin{answer}
C
\end{answer}

\begin{solution}
设动点为 \(P(u,v)\)，线段 \(P(3,0)\) 的中点为 \(M(x,y)\). 则
\(u=2x-3\)，\(v=2y\). 因为 \(P\) 在圆 \(u^2+v^2=1\) 上，所以 \((2x-3)^2+4y^2=1\).
这就是中点 \(M\) 的轨迹方程，故选 C.
\end{solution}

\begin{problem}
函数 \(y=\frac{\sin x}{|\sin x|}+\frac{|\cos x|}{\cos x}+\frac{\tan x}{|\tan x|}+\frac{|\cot x|}{\cot x}\) 的值域是（\quad）

\choices
    {\(\{-2,4\}\)}
    {\(\{-2,0,4\}\)}
    {\(\{-2,0,2,4\}\)}
    {\(\{-4,-2,0,4\}\)}
\end{problem}

\begin{answer}
B
\end{answer}

\begin{solution}
函数有意义时必须有 \(\sin x\ne0\)、\(\cos x\ne0\)，即 \(x\ne\dfrac{k\pi}{2}\)，其中 \(k\in\Z\). 在第一象限，正弦、余弦、正切、余切均为正，函数值为 \(4\)；在第二象限，正弦为正而余弦、正切、余切为负，函数值为 \(1-1-1-1=-2\)；在第三象限，正弦、余弦为负而正切、余切为正，函数值为 \(-1-1+1+1=0\)；在第四象限，只有余弦为正，函数值为 \(-1+1-1-1=-2\). 因此值域为 \(\{-2,0,4\}\)，选 B.
\end{solution}

\begin{problem}
方程 \(\sin 2x = \sin x\) 在区间 \((0,2\pi)\) 内的解的个数是（\quad）

\choices
    {\(4\)}
    {\(3\)}
    {\(2\)}
    {\(1\)}
\end{problem}

\begin{answer}
B
\end{answer}

\begin{solution}
将方程化为
\(2\sin x\cos x=\sin x\)，即 \(\sin x(2\cos x-1)=0\).
当 \(\sin x=0\) 时，在 \((0,2\pi)\) 内只有 \(x=\pi\)；当 \(2\cos x-1=0\) 时，在该区间内有 \(x=\frac{\pi}{3}\)、\(x=\frac{5\pi}{3}\). 三个解互不相同，所以解的个数为 \(3\)，选 B.
\end{solution}

\begin{problem}
函数 \(y=(x+4)^2\) 在某区间上是减函数，这区间可以是（\quad）

\choices
    {\((-\infty,-4]\)}
    {\([-4, +\infty)\)}
    {\([4, +\infty)\)}
    {\((-\infty,4]\)}
\end{problem}

\begin{answer}
A
\end{answer}

\begin{solution}
抛物线 \(y=(x+4)^2\) 的顶点横坐标为 \(-4\). 当 \(x\leqslant-4\) 时，若 \(x_1<x_2\)，则
\(x_1+4<x_2+4\leqslant0\)，
从而 \((x_1+4)^2>(x_2+4)^2\)，函数在 \((-\infty,-4]\) 上是减函数. 第二、第三个选项位于顶点右侧，函数递增；第四个选项跨过顶点，不能保持减性. 因此选 A.
\end{solution}

\begin{problem}
已知函数：
\begin{circlelist}
\item \(y=\arctan x\)，
\item \(y=\frac{\pi}{2}-\operatorname{arccot} x\).
\end{circlelist}
那么（\quad）

\choices
    {\circled{1}与\circled{2}都是奇函数}
    {\circled{1}与\circled{2}都是偶函数}
    {\circled{1}是奇函数，\circled{2}是偶函数}
    {\circled{1}与\circled{2}都不是奇函数，也都不是偶函数}
\end{problem}

\begin{answer}
A
\end{answer}

\begin{solution}
函数 \(y=\arctan x\) 的定义域为 \(\R\)，且
\(\arctan(-x)=-\arctan x\)，
所以它是奇函数.

在约定 \(\operatorname{arccot}x\in(0,\pi)\) 下，有
\(\operatorname{arccot}(-x)=\pi-\operatorname{arccot}x\).
因此对第二个函数 \(g(x)=\frac{\pi}{2}-\operatorname{arccot}x\)，有
\[
g(-x)=\frac{\pi}{2}-\operatorname{arccot}(-x)
=\frac{\pi}{2}-\left(\pi-\operatorname{arccot}x\right)
=-g(x)
\]
所以第二个函数也是奇函数，选 A.
\end{solution}

\begin{problem}
已知下图是函数 \(y=2\sin\left( \omega x+\varphi \right)\)（\(\left| \varphi \right|<\frac{\pi}{2}\)）的图象，那么（\quad）

\choices
    {\(\omega=2\)，\(\varphi=-\frac{\pi}{6}\)}
    {\(\omega=2\)，\(\varphi=\frac{\pi}{6}\)}
    {\(\omega=\frac{10}{11}\)，\(\varphi=\frac{\pi}{6}\)}
    {\(\omega=\frac{10}{11}\)，\(\varphi=-\frac{\pi}{6}\)}
\end{problem}

\begin{answer}
B
\end{answer}

\begin{solution}
原始试卷第2页的题干写有“下图”，但对应图形区域缺失，现有题面没有提供函数图象的纵截距、周期或零点位置等必要信息. 因此仅凭当前可核验的题干无法唯一确定 \(\omega\) 和 \(\varphi\)，本题的完整推导暂受题图缺失阻塞；答案表中的第二项保留为原卷记录值.
\end{solution}

\begin{problem}
如果直线 \(y=ax+2\) 与直线 \(y=3x-b\) 关于直线 \(y=x\) 对称，那么（\quad）

\choices
    {\(a=\frac{1}{3}\)，\(b=-6\)}
    {\(a=\frac{1}{3}\)，\(b=6\)}
    {\(a=3\)，\(b=-2\)}
    {\(a=3\)，\(b=6\)}
\end{problem}

\begin{answer}
B
\end{answer}

\begin{solution}
关于直线 \(y=x\) 对称会交换横、纵坐标. 直线 \(y=ax+2\) 上的点 \((x,y)\) 满足 \(y=ax+2\)，交换坐标后得到的点满足
\(x=ay+2\).
若 \(a=0\)，第一条直线为 \(y=2\)，其对称直线为 \(x=2\)，不可能与 \(y=3x-b\) 重合，因此 \(a\ne0\). 于是对称直线为 \(y=\frac{x-2}{a}\). 题设该直线就是 \(y=3x-b\)，由一次项系数和常数项分别相等，得 \(\frac{1}{a}=3\)、\(-\frac{2}{a}=-b\).
解得 \(a=\frac{1}{3}\)、\(b=6\)，因而选 B.
\end{solution}

\begin{problem}
如果抛物线 \(y^2=a(x+1)\) 的准线方程是 \(x=-3\)，那么它的焦点坐标是（\quad）

\choices
    {\((3,0)\)}
    {\((2,0)\)}
    {\((1,0)\)}
    {\((-1,0)\)}
\end{problem}

\begin{answer}
C
\end{answer}

\begin{solution}
将抛物线写成标准形式
\(y^2=4p(x+1)\).
它的顶点为 \((-1,0)\)，焦点为 \((-1+p,0)\)，准线为 \(x=-1-p\). 由题设
\(-1-p=-3\)，
得 \(p=2\). 所以焦点坐标为 \((1,0)\)，选 C.
\end{solution}

\begin{problem}
如果 \(\log_{a} 2 > \log_{b} 2 > 0 \)，那么（\quad）

\choices
    {\(1<a<b\)}
    {\(1<b<a\)}
    {\(0<a<b<1\)}
    {\(0<b<a<1\)}
\end{problem}

\begin{answer}
A
\end{answer}

\begin{solution}
对数有意义要求 \(a>0\)、\(a\ne1\)、\(b>0\)、\(b\ne1\). 因为 \(2>1\) 且
\(\log_a2>0\)、\(\log_b2>0\)，
所以 \(a>1\)、\(b>1\). 又
\(\log_a2=\frac{\ln2}{\ln a}\)、\(\log_b2=\frac{\ln2}{\ln b}\).
在 \(a\)，\(b>1\) 时，分母为正，题设 \(\log_a2>\log_b2\) 等价于 \(\ln a<\ln b\)，即 \(a<b\). 因此 \(1<a<b\)，选 A.
\end{solution}

\begin{problem}
复数 \(1+\i\cot\theta\)（\(\pi<\theta<2\pi\)）的三角形式是（\quad）

\choices
    {\(\frac{1}{\sin\theta}\left( \sin\theta+\i\cos\theta \right)\)}
    {\(\frac{1}{\sin\theta}\left[ \cos\left( \frac{\pi}{2}-\theta \right)+\i\sin\left( \frac{\pi}{2}-\theta \right) \right]\)}
    {\(-\frac{1}{\sin\theta}\left[ \cos\left( \frac{\pi}{2}+\theta \right)+\i\sin\left( \frac{\pi}{2}+\theta \right) \right]\)}
    {\(-\frac{1}{\sin\theta}\left[ \cos\left( \frac{3\pi}{2}-\theta \right)+\i\sin\left( \frac{3\pi}{2}-\theta \right) \right]\)}
\end{problem}

\begin{answer}
D
\end{answer}

\begin{solution}
因为 \(\pi<\theta<2\pi\)，所以 \(\sin\theta<0\). 先将复数化为
\(1+\i\cot\theta=\frac{\sin\theta+\i\cos\theta}{\sin\theta}\).
又
\(\cos\left(\frac{3\pi}{2}-\theta\right)=-\sin\theta\)、\(\sin\left(\frac{3\pi}{2}-\theta\right)=-\cos\theta\).
因此
\[
1+\i\cot\theta
=-\frac{1}{\sin\theta}
\left[
\cos\left(\frac{3\pi}{2}-\theta\right)
+\i\sin\left(\frac{3\pi}{2}-\theta\right)
\right]
\]
由于 \(-1/\sin\theta>0\)，这已经是三角形式，故选 D.
\end{solution}

\begin{problem}
如果 \(x\)，\(y\) 是实数，那么“\(xy>0\)”是“\(|x+y|=|x|+|y|\)”的（\quad）

\choices
    {充分但不必要条件}
    {必要但不充分条件}
    {必要而且充分的条件}
    {既不充分也不必要的条件}
\end{problem}

\begin{answer}
A
\end{answer}

\begin{solution}
对任意实数 \(x\)，\(y\)，三角不等式给出
\(|x+y|\leqslant |x|+|y|\).
等号成立当且仅当 \(x\)、\(y\) 同号或至少有一个为零，即 \(xy\geqslant0\). 因而 \(xy>0\) 时必有 \(|x+y|=|x|+|y|\)，所以它是充分条件；但例如 \(x=0\)，\(y=1\) 时等式成立而 \(xy=0\)，所以不是必要条件. 选 A.
\end{solution}

\begin{problem}
极坐标方程 \(4\rho\sin^2\frac{\theta}{2}=5\) 表示的曲线是（\quad）

\choices
    {圆}
    {椭圆}
    {抛物线}
    {双曲线的一支}
\end{problem}

\begin{answer}
C
\end{answer}

\begin{solution}
利用半角公式
\(\sin^2\frac{\theta}{2}=\frac{1-\cos\theta}{2}\)，原方程化为 \(2\rho(1-\cos\theta)=5\)，即 \(\rho(1-\cos\theta)=\frac{5}{2}\).
这是焦点在极点、偏心率为 \(1\) 的抛物线的极坐标方程，故选 C.
\end{solution}

\begin{problem}
已知等比数列的公比为 \(2\)，且前 \(4\) 项之和等于 \(1\)，那么前 \(8\) 项之和等于（\quad）

\choices
    {\(15\)}
    {\(17\)}
    {\(19\)}
    {\(21\)}
\end{problem}

\begin{answer}
B
\end{answer}

\begin{solution}
设等比数列的前 \(4\) 项和为 \(S_4\). 公比为 \(2\)，所以第 \(5\) 至第 \(8\) 项分别是第 \(1\) 至第 \(4\) 项的 \(2^4=16\) 倍，从而
\(S_8=S_4+16S_4=(1+16)\cdot1=17\).
因此选 B.
\end{solution}

\begin{problem}
函数 \(y=\tan\frac{x}{2}-\frac{1}{\sin x}\) 的最小正周期是（\quad）

\choices
    {\(\frac{\pi}{2}\)}
    {\(\pi\)}
    {\(\frac{3\pi}{2}\)}
    {\(2\pi\)}
\end{problem}

\begin{answer}
B
\end{answer}

\begin{solution}
令 \(u=\frac{x}{2}\). 在函数定义域内，\(\sin x\ne0\)，故
\[
\tan\frac{x}{2}-\frac{1}{\sin x}
=\frac{\sin u}{\cos u}-\frac{1}{2\sin u\cos u}
=\frac{2\sin^2u-1}{2\sin u\cos u}
=-\frac{\cos 2u}{\sin 2u}
=-\cot x
\]
这里函数定义域为 \(x\ne k\pi\)，\(k\in\Z\). 由于 \(\cot x\) 的基本周期为 \(\pi\)，而负号不改变周期，且 \(-\cot\left(x+\frac{\pi}{2}\right)=\tan x\ne-\cot x\)，原函数的最小正周期为 \(\pi\)，选 B.
\end{solution}

\begin{problem}
\(\tan 70^{\circ} + \tan 50^{\circ} - \sqrt{3} \tan 50^{\circ} \tan 70^{\circ}\) 的值等于（\quad）

\choices
    {\(\sqrt{3}\)}
    {\(\frac{\sqrt{3}}{3}\)}
    {\(-\frac{\sqrt{3}}{3}\)}
    {\(-\sqrt{3}\)}
\end{problem}

\begin{answer}
D
\end{answer}

\begin{solution}
令 \(u=\tan70^\circ\)、\(v=\tan50^\circ\). 由正切加法公式，
\(\tan120^\circ=\frac{u+v}{1-uv}=-\sqrt{3}\).
因此
\(u+v-\sqrt{3}uv=-\sqrt{3}\)，
正是题中所求式的值，故选 D.
\end{solution}

\begin{problem}
如果实数 \(x\)，\(y\) 满足 \((x-2)^2+y^2=3\)，那么 \(\frac{y}{x}\) 的最大值是（\quad）

\choices
    {\(\frac{1}{2}\)}
    {\(\sqrt{3}\)}
    {\(\frac{\sqrt{3}}{2}\)}
    {\(\frac{\sqrt{3}}{3}\)}
\end{problem}

\begin{answer}
B
\end{answer}

\begin{solution}
若 \(x=0\)，则由 \((x-2)^2+y^2=3\) 得 \(y^2=-1\)，矛盾，所以 \(x\ne0\). 设 \(k=\frac{y}{x}\)，直线 \(y=kx\) 与圆 \((x-2)^2+y^2=3\) 有公共点时，方程
\((x-2)^2+k^2x^2=3\)
关于 \(x\) 有实根. 整理得
\((1+k^2)x^2-4x+1=0\).
其判别式必须满足
\(\Delta=(-4)^2-4(1+k^2)=12-4k^2\geqslant0\)，
所以 \(|k|\leqslant\sqrt{3}\). 当 \(k=\sqrt{3}\) 时，方程有根 \(x=\frac12\ne0\)，确实可取到该值. 因此 \(\frac{y}{x}\) 的最大值为 \(\sqrt{3}\)，选 B.
\end{solution}

\begin{problem}
以一个正方体的顶点为顶点的四面体共有（\quad）

\choices
    {\(70\) 个}
    {\(64\) 个}
    {\(58\) 个}
    {\(52\) 个}
\end{problem}

\begin{answer}
C
\end{answer}

\begin{solution}
从正方体的 \(8\) 个顶点中任取 \(4\) 个，共有
\(\mathrm C_8^4=70\)
种取法. 将坐标平移并缩放，把正方体的顶点写成 \((0,0,0)\)、\((1,0,0)\)、\((0,1,0)\)、\((0,0,1)\)、\((1,1,0)\)、\((1,0,1)\)、\((0,1,1)\)、\((1,1,1)\). 设过四个顶点的平面方程为 \(ux+vy+wz=d\). 先看系数中有两个为零的情形，例如 \(v=w=0\)，则平面只能是 \(x=0\) 或 \(x=1\)，给出正方体的两个面；三个坐标分别作主系数时共得到 \(6\) 个面.

若恰有一个系数为零，例如 \(w=0\)，则四个顶点按 \((x,y)\) 的四种取值成两层，每一种满足 \(ux+vy=d\) 的 \((x,y)\) 都对应两个不同的 \(z\). 要得到四个顶点，四个数 \(0\)，\(u\)，\(v\)，\(u+v\) 中必须有两个相等. 因为 \(u\)，\(v\ne0\)，只能有 \(u=v\) 或 \(u=-v\)，分别得到 \(x+y=1\) 或 \(x-y=0\)；对三对坐标及两种符号共得到 \(6\) 个经过中心的对角平面.

若 \(u\)，\(v\)，\(w\) 均不为零，固定 \(x=0\) 时，四个顶点上的取值为
\(0\)、\(v\)、\(w\)、\(v+w\).
因为 \(v\)，\(w\ne0\)，每个取值至多出现 \(2\) 次；出现两次时，只可能是 \(v=w\) 时的取值 \(v\)，或 \(v=-w\) 时的取值 \(0\)，且这两种情形下都只有这一个取值重复. 固定 \(x=1\) 时，四个取值是在上述四个数上都加 \(u\). 若某个平面含有四个顶点，由于每层至多有两个顶点，它必须在两层各取两个顶点；于是 \(0\)，\(v\)，\(w\)，\(v+w\) 中的重复值 \(d\) 与 \(d-u\) 都必须重复，从而 \(u=0\)，矛盾. 因此全系数非零时不可能有四个顶点共面. 所以四点共面的组合恰有 \(6+6=12\) 组，其他组合均不共面并构成四面体. 所以四面体数为
\[
\mathrm C_8^4-6-6=70-12=58
\]
因此选 C.
\end{solution}

\begin{problem}
设 \(z\) 是复数，\(z + 2\) 的辐角为 \(\frac{\pi}{3}\)，\(z - 2\) 的辐角为 \(\frac{5\pi}{6}\)，那么 \(z\) 等于（\quad）

\choices
    {\(-\frac{1}{2}+\frac{\sqrt{3}}{2}\i\)}
    {\(-\frac{\sqrt{3}}{2}+\frac{1}{2}\i\)}
    {\(-\sqrt{3}+\i\)}
    {\(-1+\sqrt{3}\i\)}
\end{problem}

\begin{answer}
D
\end{answer}

\begin{solution}
	设 \(z=x+y\i\). 由两个辐角条件，存在 \(r\)，\(s>0\)，使
\(z+2=r\left(\cos\frac{\pi}{3}+\i\sin\frac{\pi}{3}\right)\)、\(z-2=s\left(\cos\frac{5\pi}{6}+\i\sin\frac{5\pi}{6}\right)\).
比较虚部，得
\(y=\frac{\sqrt{3}}{2}r=\frac12s\)，
所以 \(r=\frac{2y}{\sqrt{3}}\)、\(s=2y\)，并且 \(y>0\). 比较实部，得
\(x+2=\frac r2=\frac{y}{\sqrt{3}}\)、\(x-2=-\frac{\sqrt{3}}2s=-\sqrt{3}y\).
两式相减或联立可得
\(\frac{y}{\sqrt{3}}-2=2-\sqrt{3}y\)，
从而 \(y=\sqrt{3}\)、\(x=-1\). 因此 \(z=-1+\sqrt{3}\i\)，
选 D.
\end{solution}

\begin{problem}
设函数 \(y=\arctan x\) 的图象沿 \(x\) 轴正方向平移 \(2\) 个单位，所得到的图象为 \(C\). 又设图象 \(C_1\) 与 \(C\) 关于原点对称，那么 \(C_1\) 所对应的函数是（\quad）

\choices
    {\(y=-\arctan(x-2)\)}
    {\(y=\arctan(x-2)\)}
    {\(y=-\arctan(x+2)\)}
    {\(y=\arctan(x+2)\)}
\end{problem}

\begin{answer}
D
\end{answer}

\begin{solution}
函数 \(y=\arctan x\) 的图象沿 \(x\) 轴正方向平移 \(2\) 个单位后，方程为
\(y=\arctan(x-2)\).
若点 \((u,v)\) 在该图象 \(C\) 上，则 \(v=\arctan(u-2)\). 关于原点对称后的点为 \((x,y)=(-u,-v)\)，所以 \(u=-x\)、\(v=-y\)，从而
\(-y=\arctan(-x-2)=-\arctan(x+2)\).
因此 \(C_1\) 对应的函数为
\(y=\arctan(x+2)\)，
选 D.
\end{solution}

\section{填空题}

\begin{problem}
双曲线 \(\frac{x^2}{16} - \frac{y^2}{9} = 1\) 的准线方程是\fillinblank{}.
\end{problem}

\begin{answer}
\(x=\pm\frac{16}{5}\)
\end{answer}

\begin{solution}
双曲线的半实轴长为 \(a=4\)，半虚轴长为 \(b=3\)，所以半焦距
\(c=\sqrt{a^2+b^2}=\sqrt{16+9}=5\)，
偏心率为 \(e=\frac ca=\frac54\). 双曲线的准线方程为
\(x=\pm\frac ae=\pm\frac{4}{5/4}=\pm\frac{16}{5}\).
\end{solution}

\begin{problem}
\((x - 1) - (x - 1)^{2} + (x - 1)^{3} - (x - 1)^{4}\) 的展开式中，\(x^{2}\) 的系数等于\fillinblank{}.
\end{problem}

\begin{answer}
\(-10\)
\end{answer}

\begin{solution}
只有后三项含有 \(x^2\). 展开得
\[
-(x-1)^2=-x^2+2x-1
\]
\[
(x-1)^3=x^3-3x^2+3x-1
\]
\[
-(x-1)^4=-x^4+4x^3-6x^2+4x-1
\]
因此 \(x^2\) 的系数为
\(-1-3-6=-10\).
\end{solution}

\begin{problem}
已知 \(\{a_n\}\) 是公差不为零的等差数列，如果 \(S_n\) 是 \(\{a_n\}\) 的前 \(n\) 项和，那么 \(\displaystyle\lim_{n\to\infty}\frac{na_n}{S_n}=\)\fillinblank{}.
\end{problem}

\begin{answer}
\(2\)
\end{answer}

\begin{solution}
	设等差数列首项为 \(a_1\)，公差为 \(d\ne0\). 则
\(a_n=a_1+(n-1)d\)，\(S_n=na_1+\frac{n(n-1)}{2}d\).
当 \(n\) 足够大时，将分子、分母同除以 \(n^2\)，得到
\[
\frac{na_n}{S_n}
=\frac{\frac{a_1}{n}+\left(1-\frac1n\right)d}
{\frac{a_1}{n}+\frac{1}{2}\left(1-\frac1n\right)d}
\]
令 \(n\to\infty\)，由于 \(d\ne0\)，分母极限为 \(d/2\ne0\)，所以
\(\displaystyle\lim_{n\to\infty}\frac{na_n}{S_n}=\frac d{d/2}=2\).
\end{solution}

\begin{problem}
函数 \(y=\sin x+\cos x\) 的最大值是\fillinblank{}.
\end{problem}

\begin{answer}
\(\sqrt{2}\)
\end{answer}

\begin{solution}
利用辅助角公式，
\[
\sin x+\cos x
=\sqrt{2}\left(\frac{1}{\sqrt{2}}\sin x+\frac{1}{\sqrt{2}}\cos x\right)
=\sqrt{2}\sin\left(x+\frac{\pi}{4}\right)
\]
因为正弦函数的最大值为 \(1\)，所以
\(\sin x+\cos x\leqslant\sqrt{2}\).
当 \(x+\frac{\pi}{4}=\frac{\pi}{2}+2k\pi\) 时等号成立，故最大值为 \(\sqrt{2}\).
\end{solution}

\begin{problem}
如图，三棱柱 \(ABC-A_1B_1C_1\) 中，若 \(E\)，\(F\) 分别为 \(AB\)，\(AC\) 的中点，平面 \(EB_1C_1F\) 将三棱柱分成体积为 \(V_1\)，\(V_2\) 的两部分，那么 \(V_1:V_2=\)\fillinblank{}.

\bitmapfigure[max width=0.24\linewidth]{img/1990/guangdong_science/q30_fig1.png}
\end{problem}

\begin{answer}
\(7:5\)
\end{answer}

\begin{solution}
设三棱柱的底面积为 \(S\)，高为 \(h\)，则三棱柱体积为 \(Sh\). 取底面 \(ABC\) 内的向量 \(\overrightarrow{AB}\)、\(\overrightarrow{AC}\) 作为仿射坐标方向，并令距底面高为 \(th\) 的平行截面对应高度参数 \(t\)，其中 \(0\leqslant t\leqslant1\). 在这个截面内，用 \(A_t\)、\(B_t\)、\(C_t\) 表示三条侧棱上的对应点. 以 \(A_t\) 为原点，记截面内一点为 \(\lambda\overrightarrow{A_tB_t}+\mu\overrightarrow{A_tC_t}\)，则整个三角形 \(A_tB_tC_t\) 对应于 \(\lambda\geqslant0\)、\(\mu\geqslant0\)、\(\lambda+\mu\leqslant1\).

平面 \(EB_1C_1F\) 在高度 \(t\) 的截线连接了底边方向上参数和为 \(\frac12\)（当 \(t=0\) 时为 \(EF\)）与参数和为 \(1\)（当 \(t=1\) 时为 \(B_1C_1\)）的两条平行线，因而截线的方程为 \(\lambda+\mu=\frac{1+t}{2}\). 所以截线靠近 \(A_t\) 的部分是与 \(A_tB_tC_t\) 相似的三角形，相似比为 \(\frac{1+t}{2}\)，其面积为 \(S\left(\frac{1+t}{2}\right)^2\). 于是
\[
V_1=\int_0^h S\left(\frac{1+z/h}{2}\right)^2\mathrm dz
=Sh\int_0^1\left(\frac{1+t}{2}\right)^2\mathrm dt
=\frac{7}{12}Sh
\]
所以 \(V_2=Sh-V_1=\frac{5}{12}Sh\)，从而 \(V_1:V_2=7:5\).
\end{solution}

\section{解答题}

\begin{problem}
如图，在三棱锥 \(S-ABC\) 中，\(SA\perp\) 底面 \(ABC\)，\(AB\perp BC\)，\(DE\) 垂直平分 \(SC\)，且分别交 \(AC\)，\(SC\) 于 \(D\)，\(E\)，又 \(SA=AB=a\)，\(BC=\sqrt{2}a\).
\begin{enumerate}
\item 求证：\(SC\perp\) 面 \(BDE\)；
\item 求面 \(BDE\) 与面 \(BDC\) 所成的二面角的大小.
\end{enumerate}

\bitmapfigure[max width=0.28\linewidth]{img/1990/guangdong_science/q31_fig1.png}
\end{problem}

\begin{solution}
\begin{enumerate}
\item 因为 \(SA\perp\) 底面 \(ABC\)，所以 \(SA\perp AB\). 在直角三角形 \(SAB\) 中，\(SA=AB=a\)，所以 \(SB=\sqrt{2}a\). 又 \(BC=\sqrt{2}a\)，且 \(E\) 是 \(SC\) 的中点，所以在等腰三角形 \(SBC\) 中，\(BE\) 是底边 \(SC\) 的中线，从而 \(SC\perp BE\). 由 \(SC\perp DE\)，且 \(BE\cap DE=E\)，得 \(SC\perp\) 面 \(BDE\).

\item 由（1）知 \(SC\perp BD\). 又 \(SA\perp\) 底面 \(ABC\)，\(BD\subset\) 底面 \(ABC\)，所以 \(SA\perp BD\). 因为 \(SA\cap SC=S\)，所以 \(BD\perp\) 面 \(SAC\). 因而 \(DE\) 和 \(DC\) 分别是面 \(SAC\) 与面 \(BDE\)，面 \(BDC\) 的交线，\(\angle EDC\) 是所求二面角的平面角.

在直角三角形 \(ABC\) 中，\(AC=\sqrt{AB^2+BC^2}=\sqrt{3}a\). 又 \(SA\perp AC\)，所以 \(\tan\angle ACS=\frac{SA}{AC}=\frac{1}{\sqrt{3}}\)，即 \(\angle ACS=30^{\circ}\). 在直角三角形 \(DEC\) 中，\(D\in AC\)、\(E\in SC\)，故 \(\angle DCE=\angle ACS=30^{\circ}\)，且 \(\angle DEC=90^{\circ}\). 从而 \(\angle EDC=60^{\circ}\)，故所求二面角大小为 \(60^{\circ}\).
\end{enumerate}
\end{solution}

\begin{problem}
已知椭圆的焦点为 \(F_1(0,-1)\) 和 \(F_2(0,1)\)，直线 \(y=4\) 是椭圆的一条准线.
\begin{enumerate}
\item 求椭圆的方程；
\item 又设点 \(P\) 在这个椭圆上，且 \(\left| PF_1 \right|-\left| PF_2 \right|=1\)，求 \(\tan\angle F_1PF_2\) 的值.
\end{enumerate}
\end{problem}

\begin{solution}
\begin{enumerate}
\item 设椭圆方程为 \(\frac{x^2}{b^2}+\frac{y^2}{a^2}=1\)，其中 \(a>b>0\). 由焦点坐标知半焦距 \(c=1\)，又 \(y=4\) 是准线，所以 \(\frac{a^2}{c}=4\)，从而 \(a^2=4\)，\(b^2=a^2-c^2=3\). 因此椭圆方程为 \(\frac{y^2}{4}+\frac{x^2}{3}=1\).

\item 由椭圆定义，\(\left|PF_1\right|+\left|PF_2\right|=2a=4\). 结合题设 \(\left|PF_1\right|-\left|PF_2\right|=1\)，得 \(\left|PF_1\right|=\frac{5}{2}\)，\(\left|PF_2\right|=\frac{3}{2}\). 又 \(\left|F_1F_2\right|=2\)，在三角形 \(F_1PF_2\) 中由余弦定理，
\[
\cos\angle F_1PF_2
=\frac{\left|PF_1\right|^2+\left|PF_2\right|^2-\left|F_1F_2\right|^2}{2\left|PF_1\right|\left|PF_2\right|}
=\frac{\frac{25}{4}+\frac{9}{4}-4}{\frac{15}{2}}
=\frac{3}{5}
\]
所以
\[
\tan\angle F_1PF_2
 =\sqrt{\frac{1}{\cos^2\angle F_1PF_2}-1}
 =\sqrt{\frac{25}{9}-1}
 =\frac{4}{3}
\]
\end{enumerate}
\end{solution}

\begin{problem}
已知数列 \(\{a_n\}\) 的前 \(n\) 项和的公式是 \(S_n=\frac{\pi}{12}(2n^2+n)\).
\begin{enumerate}
\item 求证 \(\{a_n\}\) 是等差数列，并求出它的首项和公差；
\item 记 \(b_n=\sin a_n\sin a_{n+1}\sin a_{n+2}\)，求证：对任何自然数 \(n\)，都有 \(b_n=\frac{\sqrt{2}}{8}(-1)^{n-1}\).
\end{enumerate}
\end{problem}

\begin{solution}
\begin{enumerate}
\item 因为 \(\{a_n\}\) 的前 \(n\) 项和为 \(S_n=\frac{\pi}{12}(2n^2+n)\)，所以 \(a_1=S_1=\frac{\pi}{4}\). 当 \(n\geqslant2\) 时，\(a_n=S_n-S_{n-1}=\frac{\pi}{12}(4n-1)\)，该式对 \(n=1\) 也成立. 又 \(a_{n+1}-a_n=\frac{\pi}{3}\)，所以 \(\{a_n\}\) 是首项为 \(\frac{\pi}{4}\)，公差为 \(\frac{\pi}{3}\) 的等差数列.

\item 由（1）得 \(a_n=\frac{\pi}{12}(4n-1)\). 因为对任意自然数 \(n\)，\(4n-1\) 是奇数，所以 \(\sin a_n\neq0\). 先算 \(b_1\)：
\[
b_1=\sin\frac{\pi}{4}\sin\frac{7\pi}{12}\sin\frac{11\pi}{12}
=\frac{\sqrt{2}}{2}\cdot\frac{1}{2}\left(\cos\frac{\pi}{3}-\cos\frac{3\pi}{2}\right)
=\frac{\sqrt{2}}{8}
\]
又因为 \(\{a_n\}\) 的公差为 \(\frac{\pi}{3}\)，所以 \(a_{n+3}=a_n+\pi\)，从而 \(\displaystyle\frac{b_{n+1}}{b_n}=\frac{\sin a_{n+3}}{\sin a_n}=-1\). 因此 \(\{b_n\}\) 是首项为 \(\frac{\sqrt{2}}{8}\)，公比为 \(-1\) 的等比数列，故对任意自然数 \(n\)，\(\displaystyle b_n=\frac{\sqrt{2}}{8}(-1)^{n-1}\).
\end{enumerate}
\end{solution}

\begin{problem}
设 \(f(x)=\lg\frac{1+2^x+4^xa}{3}\). 其中 \(a\in\R\).
\begin{enumerate}
\item 如果 \(0<a\leqslant1\)，求证：当 \(x\neq0\) 时，有 \(2f(x)<f(2x)\)；
\item 如果当 \(x\in(-\infty,1]\) 时 \(f(x)\) 有意义，求 \(a\) 的取值范围.
\end{enumerate}
\end{problem}

\begin{solution}
\begin{enumerate}
\item 记 \(p=1\)，\(q=2^x\)，\(r=4^xa\). 因为 \(0<a\leqslant1\)，所以 \(p\)、\(q\)、\(r\) 均为正数；又因为 \(x\neq0\)，所以 \(p\neq q\)，从而 \(p^2+q^2>2pq\)，并且 \(q^2+r^2\geqslant2qr\)、\(r^2+p^2\geqslant2rp\). 由三式相加可得
\[
\frac{1}{9}(p+q+r)^2
<\frac{1}{9}\left(3p^2+3q^2+3r^2\right)
=\frac{1}{3}(p^2+q^2+r^2)
=\frac{1}{3}\left(1+2^{2x}+4^{2x}a^2\right)
\leqslant\frac{1}{3}\left(1+2^{2x}+4^{2x}a\right)
\]
另一方面，\(2f(x)=\lg\frac{(p+q+r)^2}{9}\)，而 \(f(2x)=\lg\frac{1+2^{2x}+4^{2x}a}{3}\). 上面的严格不等式以及常用对数函数的单调性给出 \(2f(x)<f(2x)\).

\item 当 \(x\in(-\infty,1]\) 时，\(f(x)\) 有意义的充要条件是 \(1+2^x+4^xa>0\). 因为 \(4^x>0\)，这等价于 \(\displaystyle a>-\left[\left(\frac{1}{4}\right)^x+\left(\frac{1}{2}\right)^x\right]\). 函数 \(-\left(\frac{1}{4}\right)^x-\left(\frac{1}{2}\right)^x\) 在 \((-\infty,1]\) 上递增，在 \(x=1\) 时取得最大值 \(-\frac{3}{4}\). 因而所求取值范围为 \(\displaystyle a>-\frac{3}{4}\).
\end{enumerate}
\end{solution}
